\subsection{例子：德拉姆复形}

在本节中，我们假设 A 是交换环 \(k\) 上的一个交换 \(k\) -代数。用 \(\Omega_{\mathbf{A}/k}^{1}\) 表示 A 的凯勒微分 A-模（III，第 134 页），用 \(d^{0}: \mathbf{A} \to \Omega_{\mathbf{A}/k}^{1}\) 表示 \(k\) -导子 \(d_{\mathbf{A}/k}\) ，并用 \(\Omega_{\mathbf{A}/k}\) 表示分次 \(k\) -代数 \(\boldsymbol{\Lambda}_{\mathbf{A}}(\Omega_{\mathbf{A}/k}^{1})\) .

命题13. --- 存在唯一的次数为 1、平方为零的 k-反导子 \(d: \Omega_{\mathbf{A}/k} \to \Omega_{\mathbf{A}/k}\) ，它延拓导子 \(d^{0}: \mathbf{A} \to \Omega_{\mathbf{A}/k}^{1}\) .

让我们证明反导子 \(d\) 的唯一性。由于 \(d \circ d = 0\) ，我们有，对于 \(y, x_1, \ldots, x_p \in \mathbf{A}\) :

\[
d \left(y d x _ {1} \wedge \dots \wedge d x _ {p}\right) = d y \wedge d x _ {1} \wedge \dots \wedge d x _ {p}.
\]

A-模 \(\Omega_{\mathbf{A}/k}^{p}\) 由元素 \(dx_{1} \wedge \ldots \wedge dx_{p}\) 生成，这就证明了 \(d\) 的唯一性 .

为了证明存在性，只需构造一个 \(k\) -同态 \(d^{1}:\Omega_{\mathbb{A}/k}^{1}\to\Omega_{\mathbb{A}/k}^{2}\) ，使得 \(d^{1}\circ d^{0}=0\) 且

\[
d ^ {1} (a \omega) = d ^ {0} (a) \wedge \omega + a d ^ {1} (\omega) \quad \text {for} \quad a \in A, \quad \omega \in \Omega_ {A k} ^ {1}.\tag{11}
\]

事实上，由 III，第 128 页，命题 14（考虑到 III，第 118 页，注记 2）可知，存在一个反导子 \(d: \Omega_{\mathbf{A}/k} \to \Omega_{\mathbf{A}/k}\) ，它在次数为 0 时与 \(d^{0}\) 重合，在次数为 1 时与 \(d^{1}\) 重合。由于 \(d^{0}\) 在 \(\mathbf{A}\) 上为零，反导子 \(d\) 是 \(k\) -线性的；由于 \(d^{1} \circ d^{0} = 0\) ，我们有 \(d \circ d = 0\) ，因为 \(\Omega_{\mathbf{A}/k}\) 作为 \(\mathbf{A}\) -代数由元素 \(d^{0} a\) （ \(a \in \mathbf{A}\) ）生成 .

为了定义 \(d^{1}\) ，回顾（III，第 133 页） \(\Omega_{\mathbf{A}/k}^{1}\) 等于 A-模 \(\mathfrak{I}/\mathfrak{I}^{2}\) ，其中 \(\mathfrak{I}\) 是乘法 \(m: \mathbf{A} \otimes_{k} \mathbf{A} \to \mathbf{A}\) 的核。考虑 \(k\) -线性映射 \(u: \mathbf{A} \otimes_{k} \mathbf{A} \to \Omega_{\mathbf{A}/k}^{2}\) ，它由 \(u(x \otimes y) = d^{0}(y) \wedge d^{0}(x)\) 定义。我们有

\[
u (a x \otimes y - x \otimes a y) = d ^ {0} (y) \wedge d ^ {0} (a x) - d ^ {0} (a y) \wedge d ^ {0} (x) = d ^ {0} (x y) \wedge d ^ {0} (a)
\]

对于 A 中的 x、y 和 a，由此可得

\[
u ((a \otimes 1 - 1 \otimes a) \xi) = d ^ {0} (m (\xi)) \wedge d ^ {0} (a), \quad \xi \in \mathrm{A} \otimes_ {k} \mathrm{A}, \quad a \in \mathrm{A}. \tag {12}
\]

由于 \(\mathfrak{I}\) 作为左 A-模由元素 ( \(a \otimes 1 - 1 \otimes a\) ) （ \(a \in \mathbf{A}\) ）生成，我们推出 \(u(\mathfrak{I}^2) = 0\) ；因此 \(u\) 通过限制到 \(\mathfrak{I}\) 并过渡到商，定义了一个 \(k\) -线性映射 \(d^1: \mathfrak{I}/\mathfrak{I}^2 \to \Omega_{\mathrm{A}/k}^2\) .

在（12）中取 \(\xi = b \otimes 1\) （其中 \(b \in \mathbf{A}\) ），我们得到 \(d^{1}(bd^{0}(a)) = d^{0}(b) \wedge d^{0}(a)\) ；由此得出 \(d^{1} \circ d^{0} = 0\) ，以及 \(d^{1}(c\omega) = d^{0}(c) \wedge \omega + cd^{1}(\omega)\) （对 \(c \in \mathbf{A}\) 与 \(\omega = bd^{0}(a)\) ）。由于 \(\Omega_{\mathbf{A}/k}^{1}\) 作为 \(k\) -模由元素 \(bd^{0}(a)\) （ \(a\) 与 \(b\) 属于 \(\mathbf{A}\) ）生成，公式（11）对所有 \(\omega \in \Omega_{\mathbf{A}/k}\) 成立，这就完成了命题的证明。

有时称元素 \(\omega\in\Omega_{A/k}^{p}\) 为次数为 \(p\) 的外微分形式（A 在 \(k\) 上），并称反导子 \(d\) 为 \(\Omega_{A/k}\) 的外微分；复形（ \(\Omega_{A/k}\) , \(d\) ）称为 A 在 \(k\) 上的德拉姆复形，其同调为 A 在 \(k\) 上的德拉姆上同调 .

例 1. --- 取 A 为环 \(k[\mathbf{X}_{1},\ldots,\mathbf{X}_{n}]\) 。则 \(\Omega_{\mathbf{A}/k}^{1}\) 是一个自由 A-模，其基为 \(d\mathbf{X}_{1},\ldots,d\mathbf{X}_{n}\) （III，第 134 页，例）。因此，若对每个子集 \(\mathrm{I}=\{i_{1},\ldots,i_{p}\}\) （ \([1,n]\) 的）设 \(d\mathbf{X}_{\mathrm{I}}=d\mathbf{X}_{i_{1}}\wedge\ldots\wedge d\mathbf{X}_{i_{p}}\) （其中 \(i_{1}<\ldots<i_{p}\) ），则 A-模 \(\Omega_{\mathbf{A}/k}^{p}\) 以元素 \(d\mathbf{X}_{\mathrm{I}}\) 为基，其中 I 遍历 \([1,n]\) 的基数为 \(p\) 的子集的集合。我们有

\[
d (\mathrm{P} d \mathrm{X} _ {\mathrm{I}}) = d \mathrm{P} \wedge d \mathrm{X} _ {\mathrm{I}} = \sum_ {i \notin \mathrm{I}} (- 1) ^ {n (\mathrm{I}, i)} \frac {\partial \mathrm{P}}{\partial \mathrm{X} _ {i}} d \mathrm{X} _ {\mathrm{I} \cup \{i \}},
\]

其中 \(n(\mathbf{I}, i)\) 表示 I 中严格小于 \(i\) 的元素的个数 .

\(Z^{p}(\Omega_{\mathbf{A}/k})\) 的闭链因此是这样的元素 \(\omega = \sum_{\text{Card (I)} = p} \mathrm{P}_{\mathrm{I}} d\mathrm{X}_{\mathrm{I}}\) ：对每个子集 \(J\) （含 \((p + 1)\) 个元素，属于 \([1, n]\) ），有：

\[
\sum_ {i \in J} (- 1) ^ {n (J, i)} \frac {\partial P _ {J - \{i \}}}{\partial X _ {i}} = 0.
\]

元素 ω 是一个边缘，如果对于每个子集 \(J \subset [1, n]\) （含 \((p - 1)\) 个元素），可以选择一个多项式 \(Q_{J} \in A\) ，使得：

\[
\mathrm{P} _ {\mathrm{I}} = \sum_ {j \in \mathrm{I}} (- 1) ^ {n (\mathrm{I}, j)} \frac {\partial \mathrm{Q} _ {\mathrm{I} - \{j \}}}{\partial \mathrm{X} _ {j}}.
\]

我们将在 § 9 中看到，如果 k 是一个 Q-代数，则 A 在 k 上的德拉姆复形在次数 \textgreater{} 0 上是零调的（X，第 159 页，注记 4）。

* 例 2. --- 假设 \(k = \mathbf{C}\) 且 \(\mathbf{A} = \mathbf{C}\) \([\mathrm{X}_1, \ldots, \mathrm{X}_n] / (\mathrm{P}_1, \ldots, \mathrm{P}_r)\) ，其中 \(\mathrm{P}_i\) 是 \(\mathrm{X}_1, \ldots, \mathrm{X}_n\) 的多项式，使得 \(\mathbf{C}^n\) 中所有 \(\mathrm{P}_i\) 为零的点集是 \(\mathbf{C}^n\) 的一个解析子流形 V。可以证明 A 在 C 上的德拉姆上拟同构于奇异上同调 H(V, C)。*

设 M 是一个 A-模，且 \(\nabla^{0}\) 为 \(k\) -线性映射，从 M 到 \(\mathbf{M} \otimes_{\mathbf{A}} \Omega_{\mathbf{A}/k}^{1}\) 使得

\[
\nabla^ {0} (a m) = a \nabla^ {0} (m) + m \otimes d a \quad \text { for } \quad a \in \mathrm{A}, \quad m \in \mathrm{M}\tag{13}
\]

（有时也称 \(\nabla^{0}\) 为 A-模 M 上的一个联络）。

命题 14. --- (i) 存在唯一的 k-线性映射 \(\nabla\) 从 \(\Omega_{\mathbf{A}/k}\) 右模 \(\mathbf{M} \otimes_{\mathbf{A}} \Omega_{\mathbf{A}/k}\) 到自身，次数为 1 的分次映射，它扩展了 \(\nabla^{0}\) 在次数0中，并满足恒等式：

\[
\nabla (x \omega) = (\nabla x) \omega + (- 1) ^ {p} x (d \omega) \quad \mathrm{for} \quad x \in \mathrm{M} \otimes_ {\mathrm{A}} \Omega_ {\mathrm{A} / k} ^ {p}, \quad \omega \in \Omega_ {\mathrm{A} / k}. \tag {14}
\]

\begin{enumerate}
\def\labelenumi{(\roman{enumi})}
\setcounter{enumi}{1}
\tightlist
\item
  复合映射 \(\nabla \circ \nabla\) 是 \(\Omega_{\mathbf{A} / k}\) -线性；特别地，映射 \(\mathbf{R} = \nabla^{1} \circ \nabla^{0}\) 的 \(\mathbf{M}\) 到 \(\mathbf{M} \otimes_{\mathbf{A}} \Omega_{\mathbf{A} / k}^{2}\) 是 \(\mathbf{A}\) -线性，并且我们有
\end{enumerate}

\[
\nabla \circ \nabla (m \otimes \omega) = R (m). \omega \quad \mathrm{for} \quad m \in M, \quad \omega \in \Omega_ {A / k}.
\]

同态 R 有时称为联络 \(\nabla^{0}\) 的曲率同态；若它为零，则偶对 \((\mathbf{M} \otimes_{\mathbf{A}} \Omega_{\mathbf{A}/k}, \nabla)\) 是一个复形，也称为 \((\mathbf{M}, \nabla^{0})\) 在 \(k\) 上的德拉姆复形 .

我们来证明 (i)。 \(\nabla\) 的唯一性是显然的。让我们定义一个 \(k\) -同态 \(\overline{\nabla}\) ，从 \(\mathbf{M} \otimes_k \Omega_{\mathbf{A}/k}\) 到 \(\mathbf{M} \otimes_{\mathbf{A}} \Omega_{\mathbf{A}/k}\) ：

\[
\overline {{{\nabla}}} (m \otimes_ {k} \omega) = (\nabla^ {0} m) \omega + m \otimes d \omega \quad \text { for } \quad m \in \mathbf {M}, \quad \omega \in \Omega_ {\mathrm{A} / k}.
\]

由（13）可得 \(\overline{\nabla}(am \otimes \omega) = \overline{\nabla}(m \otimes a\omega)\) ，因此通过过渡到商，我们得到一个 \(k\) -同态 \(\nabla\) ，它从 \(\mathbf{M} \otimes_{\mathbf{A}} \Omega_{\mathbf{A}/k}\) 到自身、次数为 1 的分次映射，且在次数 0 时延拓 \(\nabla^{0}\) 。让我们验证（14）：我们有，对于 \(m \in \mathbf{M}\) , \(\alpha \in \Omega_{\mathbf{A}/k}^{p}\) , \(\omega \in \Omega_{\mathbf{A}/k}\) :

\[
\begin{array}{r l} \nabla ((m \otimes \alpha). \omega) & = \nabla (m \otimes (\alpha \wedge \omega)) = \nabla^ {0} (m). (\alpha \wedge \omega) + m \otimes d (\alpha \wedge \omega) \\ & = \nabla^ {0} (m) \alpha . \omega + (m \otimes d \alpha) \omega + (- 1) ^ {p} (m \otimes \alpha) d \omega \\ & = (\nabla (m \otimes \alpha)) \omega + (- 1) ^ {p} (m \otimes \alpha) d \omega \end{array}
\]

这证明了 (14) 式对 \(x = m \otimes \alpha\) 成立；一般情形由此通过线性性推出。

让我们证明(ii)。设 \(x \in \mathbf{M} \otimes_{\mathbf{A}} \Omega_{\mathbf{A}/k}^{p}\) , \(\omega \in \Omega_{\mathbf{A}/k}\) ；通过反复应用(14)，我们得到：

\[
\begin{array}{r l} \nabla \circ \nabla (x \omega) = & \nabla (\nabla (x) \omega) + (- 1) ^ {p} \nabla (x d \omega) \\ & = (\nabla \circ \nabla (x)) \omega + (- 1) ^ {p + 1} \nabla (x) (d \omega) + (- 1) ^ {p} \nabla (x) (d \omega) \\ & = (\nabla \circ \nabla (x)) \omega , \end{array}
\]

这证明了(ii)的第一个断言；其余断言随之立即成立。

